\chapter{Reinforcement Learning}

Reinforcement learning differs from supervised learning in one crucial way: the learner does not receive the correct target for each action immediately. It must act, observe consequences, and infer which decisions were responsible for later reward.

\section{Markov decision processes}

The standard formal model is a Markov decision process (MDP) with states $s_t$, actions $a_t$, rewards $r_t$, transition law $p(s_{t+1}\mid s_t,a_t)$, and policy $\pi(a\mid s)$.

\begin{figure}[t]
\centering
\begin{tikzpicture}[node distance=2.6cm, >=Latex]
  \node[draw, rounded corners, minimum width=3.2cm, minimum height=1.2cm] (env) {environment};
  \node[draw, rounded corners, minimum width=3.2cm, minimum height=1.2cm, right=of env] (agent) {agent};
  \draw[->, very thick] (env) -- node[above] {$s_t, r_t$} (agent);
  \draw[->, very thick] (agent) to[bend left=15] node[below] {$a_t$} (env);
\end{tikzpicture}
\caption{The reinforcement-learning loop: action changes future state distributions, so data are generated by the current policy itself.}
\end{figure}

The return from time $t$ is
\[
  G_t = \sum_{k=0}^{\infty} \gamma^k r_{t+k+1},
  \qquad 0 \le \gamma < 1.
\]
Discounting makes distant reward less important and ensures convergence of the infinite-horizon sum under mild conditions.

\section{Value functions}

For policy $\pi$, the state-value and action-value functions are
\[
  V^\pi(s) = \E_\pi[G_t \mid s_t=s],
  \qquad
  Q^\pi(s,a) = \E_\pi[G_t \mid s_t=s, a_t=a].
\]
These satisfy Bellman equations such as
\[
  V^\pi(s)
  =
  \sum_a \pi(a \mid s)
  \sum_{s',r}
  p(s',r \mid s,a)
  \bigl[r + \gamma V^\pi(s')\bigr].
\]

\begin{keyequation}[Bellman recursion]
\[
V^\pi(s)
=
\sum_a \pi(a \mid s)
\sum_{s',r}
p(s',r \mid s,a)
\bigl[r + \gamma V^\pi(s')\bigr]
\]
\tcblower
The value of a state is the expected immediate reward plus the discounted value of the next state. Reinforcement learning works because long-horizon evaluation can be written recursively.
\end{keyequation}

\section{Temporal-difference learning}

Temporal-difference (TD) learning replaces the unknown future return by a one-step bootstrap target:
\[
  V(s_t)
  \leftarrow
  V(s_t) + \alpha
  \bigl[
    r_{t+1} + \gamma V(s_{t+1}) - V(s_t)
  \bigr].
\]
The bracketed term is the TD error. It measures how surprising the next transition was relative to the current value estimate.

\section{Control: SARSA and Q-learning}

For action values, two central updates are:
\[
  Q(s_t,a_t)
  \leftarrow
  Q(s_t,a_t)
  +
  \alpha
  \bigl[
    r_{t+1}
    +
    \gamma Q(s_{t+1},a_{t+1})
    -
    Q(s_t,a_t)
  \bigr]
\]
for SARSA, and
\[
  Q(s_t,a_t)
  \leftarrow
  Q(s_t,a_t)
  +
  \alpha
  \bigl[
    r_{t+1}
    +
    \gamma \max_{a'} Q(s_{t+1},a')
    -
    Q(s_t,a_t)
  \bigr]
\]
for Q-learning.

SARSA is on-policy: it learns the value of the behavior actually being followed. Q-learning is off-policy: it learns toward a greedy target even while behavior remains exploratory.

\section{Exploration and exploitation}

An agent that only exploits current knowledge may never discover better actions. An agent that only explores never consolidates what it has learned. Methods such as $\varepsilon$-greedy policies make this trade-off explicit:
\begin{itemize}
  \item with probability $1-\varepsilon$, choose the current best action,
  \item with probability $\varepsilon$, explore.
\end{itemize}

\begin{aside}[Why reinforcement learning is hard]
The data distribution changes as the policy changes. In supervised learning, the dataset is usually fixed before optimization begins. In reinforcement learning, the learner partly creates its own future training data.
\end{aside}

\section{Course synthesis}

The course closes on a useful unification:
\begin{itemize}
  \item supervised learning studies prediction from examples,
  \item statistical learning theory studies when such prediction generalizes,
  \item Bayesian methods attach uncertainty to the predictions,
  \item reinforcement learning adds action and delayed credit assignment.
\end{itemize}

The common thread is always the same: choose a representation, define a principled objective, and use mathematics to separate what is merely fit from what is genuinely learned.
